\introduction

% ЧЕРНОВИК
Совместное обучение моделей на данных нескольких организаций ограничено
требованиями к конфиденциальности: передать исходные данные третьей стороне
часто невозможно. Технологии многосторонних конфиденциальных вычислений (МКВ)
позволяют выполнять обучение и вывод над данными в виде долей секрета, не
раскрывая их отдельным участникам~\cite{zapechnikov2021two,zapechnikov2022three}.

В НИРС~1 были рассмотрены угрозы безопасности систем машинного
обучения~\cite{ilyushin2022attacks,papernot2018sok}. Настоящая работа
продолжает её: определяется, какие из этих угроз закрывают МКВ, и
проектируется система на их основе.

\textbf{Цель работы:} спроектировать защищённую систему машинного обучения с
применением МКВ и подготовить экспериментальную проверку.

\textbf{Задачи:}
\begin{enumerate}
    \item сопоставить угрозы из НИРС~1 с возможностями МКВ;
    \item построить детальную модель нарушителя и угроз;
    \item разработать архитектуру системы и требования к ней;
    \item обосновать выбор протокола ABY3 и платформы TF Encrypted, параметров
    и способа вычисления решающего дерева;
    \item разработать методику экспериментов;
    \item выполнить пробный запуск на одной машине.
\end{enumerate}
